%%%%%%%%%%%%%%%%%%%%%%%%%%%%%%%%%%%%%
%%% !!! DO NOT EDIT THIS FILE !!! %%%
%%%%%%%%%%%%%%%%%%%%%%%%%%%%%%%%%%%%%
\documentclass{article}

\usepackage{geometry}
\usepackage{fontspec}
\usepackage[svgnames]{xcolor}
\usepackage{hyperref}
\usepackage{paracol}
\usepackage{ragged2e}
\usepackage{amsmath, amssymb, physics}

\geometry{paperheight=29.7cm, paperwidth=21cm, includehead, includefoot, nomarginpar, textheight = 26cm, textwidth=20cm, headheight=1cm}
\setmainfont[AutoFakeBold=6.5, AutoFakeSlant]{Montserrat Regular}
\hypersetup{colorlinks=true, urlcolor=blue}
\columnratio{0.75}

\newcommand{\abstracttitle}[1]{\def\titleVar{#1}}

\newcommand{\speaker}[2]{\underline{#1}$^{#2}$}
\newcommand{\coauthor}[2]{#1$^{#2}$}
\newcommand{\authorslist}[1]{\def\authorsVar{#1}}

\newcounter{affilnum}
\newcommand{\affiliation}[1]{
    \addtocounter{affilnum}{1}$^{\arabic{affilnum}}$#1\\
}
\newcommand{\affiliations}[1]{
    \def\affilationsVar{#1}
    \setcounter{affilnum}{0}
}

\newcommand{\keywords}[1]{\def\keywordsVar{#1}}
\renewcommand{\abstract}[1]{\def\abstractVar{#1}}

\def\changemargin#1#2{\list{}{\rightmargin#2\leftmargin#1}\item[]}
\let\endchangemargin=\endlist 

\newcommand{\createabstract}[1]{
\input{#1}

\begin{changemargin}{10pt}{10pt} 
\begin{paracol}{2}
\hspace{-10pt}\textbf{\titleVar}
\vspace{4pt}\\
\textcolor{Tomato}{\authorsVar}
\vspace{4pt}\\
\textcolor{Teal}{\affilationsVar}
\switchcolumn
\textbf{Keywords:}
\keywordsVar
\vspace{4pt}\\
% \textbf{Poster ID:} XXX
\end{paracol}
\vspace{-0.5cm}
\abstractVar
\end{changemargin}
}

\title{Abstract Template
\\
METANANO School 2024 On
\\
Nanophotonics and Advanced Materials}
\author{}
\date{\vspace{-1.5cm}}

\begin{document}
\maketitle
\subsection*{How to use this template}
\begin{enumerate}
    \item If you don't know how to use \LaTeX{}, you can read \href{https://www.overleaf.com/learn/latex/Learn_LaTeX_in_30_minutes}{this guide}.
    If you don't know how to use Overleaf, you can read 
    \href{https://www.overleaf.com/learn}{{this guide}}.
    
    \item On your left you can find Overleaf file tree.
    If not, click on the small arrow located at the center of the left border of your screen to open it. Here you can find two documents: "main.tex" and "abstract\_info.tex".
    \textbf{Edit only "abstract\_info.tex"}, do not touch "main.tex".

    \item In "abstract\_info.tex" you will need to fill out following fields: poster title, authors list, affiliations, keywords, and abstract.
    \begin{enumerate}
        \item \textbf{Abstract title}
        \\
        Write here your poster title.
        Capitalize "major" words of your titile (nouns, pronouns, verbs, adjectives, adverbs, and some conjunctions).
        Lowercase the conjunctions "and", "but", "for", "or", and "nor", and the articles "the", "a", and "an".
        Usage of inline math mode is allowed for chemical formulas, but please, do not write mathematical equations in the title, and try to avoid Greek letters and math mode in title in general.
        \\
        For example:
        \begin{changemargin}{2cm}{0cm}
        \begin{FlushLeft}
            \verb|\postertitle{|
            \\
            \begin{changemargin}{0.7cm}{0cm}
            \verb|Engineering Gamma and Alpha States In $A_2 B_4$ Systems|
            \end{changemargin}
            \verb|}|
        \end{FlushLeft}
        \end{changemargin}
        will produce following title
        \abstracttitle{
            Engineering Gamma and Alpha States In $A_2 B_4$ Systems
        }
        \begin{center}
            \textbf{\titleVar}
        \end{center}

        \item \textbf{Authors list}
        \\
        Write all the authors of your work here.
        Use \verb|\speaker{}{}| command to highlight the speaker and
        \\
        \verb|\coauthor{}{}| command to list all other authors.
        The first braces are for the name, second are for the affiliation number.
        List all authors in the order you want them to be in the book of abstracts.
        Separate authors by commas.
        \\
        For example:
        \begin{changemargin}{2cm}{0cm}
        \begin{FlushLeft}
            \verb|\authorslist{|
            \\
            \begin{changemargin}{0.7cm}{0cm}
            \verb|\coauthor{Keith Moon}{1},|
            \\
            \verb|\speaker{Rogel Daltrey}{1, 2},| 
            \\
            \verb|\coauthor{John Entwistle}{1},|
            \\
            \verb|\coauthor{Pete Townshend}{1,2}.|
            \\
            \end{changemargin}
            \verb|}|
        \end{FlushLeft}
        \end{changemargin}
        will produce following authors list
        \authorslist{
            \coauthor{Keith Moon}{1},
            \speaker{Roger Daltery}{1,2},
            \coauthor{John Entwistle}{1},
            \coauthor{Pete Townshend}{1,2}.
        }
        \begin{center}
            \textcolor{Tomato}{\authorsVar}
        \end{center}

        \item \textbf{Affiliations}
        \\
        List here affiliations starting from the first.
        Use \verb|\affiliation{}| command for each affiliation.
        Please, try to maintain the following affiliation structure: "Lab/Faculty, University, City, Country, Zip Code".
        \\
        For example:
        \begin{changemargin}{2cm}{0cm}
        \begin{FlushLeft}
            \verb|\affiliations{|
            \\
            \begin{changemargin}{0.7cm}{0cm}
            \verb|\affiliation{Physics Lab, University of Science, Some Town, Some Country, 069420}|
            \\
            \verb|\affiliation{Disco Lab, University of Party, Another Town, Another Country, 080085}|
            \end{changemargin}
            \verb|}|
        \end{FlushLeft}
        \end{changemargin}
        will produce following affiliations
        \affiliations{
            \hspace{-0.2cm}\affiliation{Physics Lab, University of Science, Some Town, Some Country, 069420}
            \affiliation{Disco Lab, University of Party, Another Town, Another Country, 080085}
        }
        \begin{center}
            \textcolor{Teal}{\affilationsVar}
        \end{center}

        \item \textbf{Keywords}
        \\
        List here your keywords.
        Please, write at least 2 but no more than 6 keywords.
        Long-tail keywords like "topological insulators" or "bound states in the continuum" are allowed, but, if possible, try to use acronyms, e.g. "BIC" for "bound states in the continuum", or EFISH for "electric-field induced second harmonic".

        \item \textbf{Abstract}
        \\
        Write here your abstract and make sure that it compiles successfully.
        There is no limit to the number of words in the abstract, but it is recommended to write at least 50 words, but no more than 250 words.
        If you need, you can use inline equations, for example: $f(x) = x^2 - c$.
        If you need to write fractions, write them in one row, like this: $\alpha = \beta/2\gamma$.
        Please, do not use any additional packages.
        \textbf{And, please, do not include figures.}
    \end{enumerate}

    \item If it will be unclear for you how to fill in some part of the "abstratc\_info.tex", you can consult "abstract\_info.tex" in \textbf{filled in example of template}, which is available via the \href{https://www.overleaf.com/read/vwgjhqtdrhqd#601482}{following link}.

    \item After you wrote all the information, please, compile the document (\href{https://www.overleaf.com/learn/how-to/How_do_I_compile_my_tex_file%3F}{guide here}).
    If Overleaf is unable to create PDF due to errors, please, correct them and recompile the document.
    If the document is compiled successfully, below you will be able to see how your abstract will look in the book of abstracts.
    Please, check everything for errors and typos.
    Also, check the correctness of the title of your poster, names of the authors, affiliations, and other information.
    Please note that the information you wrote will not be checked and changed by the organizers.
    It will be presented in the book of abstracts in the way it is presented by you in the template.  

    \item If everything is compiled successfully, if there are no typos and errors, and if you are happy with the way your abstract looks, please, download your project as ".zip" file (\href{https://www.overleaf.com/learn/how-to/Exporting_your_work_from_Overleaf#Downloading_Overleaf_projects}{guide here}), rename it as "YourName\_YourSurname.zip".
    
\end{enumerate}

\subsection*{This Is How Your Abstract Will Look}
\vspace{-0.4cm}\noindent\makebox[\linewidth]{\rule{\paperwidth}{0.1pt}}\vspace{-0.4cm}
\createabstract{abstract_info}
\noindent\makebox[\linewidth]{\rule{\paperwidth}{0.1pt}}

\end{document}
